
\def\PUDcodigo{ECA226}

\def\PUDcodacad{04507.44}

\def\PUDdisciplina{Robótica I}

\def\PUDsemestre{S9}

\def\PUDhoras{80}

%NOVOS ---------------------------------------------------
\def\PUDhorasTeoria{40}
\def\PUDhorasPratica{40}

\def\PUDinterECA{ 
\hyperref[PUD:ECA225]{Dispositivos Periféricos (04507.42)} 

\hyperref[PUD:ECA227]{SDCD (04507.45)} 
}

\def\PUDatuaisECA{- Aplicações com tecnologias atuais}

\def\PUDrecursos{Projetor multimídia; Quadro branco e pincel; Aulas em laboratório.}

%----------------------------------------------------------

\def\PUDcreditos{4}

\def\PUDprerequisito{ \hyperref[PUD:ECA225]{ECA225} }

\def\PUDpreacad{ \hyperref[PUD:ECA225]{Dispositivos Periféricos (04507.42)} }

\def\PUDobjetivos{Compreender as técnicas básicas de modelagem e programação de robôs industriais tanto manipuladores articulados quanto móveis. Ler, interpretar e escrever programasna linguagem de robôs industriais. Além de verificar o contexto de aplicações regionais e interdiciplinariedade tanto com as disciplinas de pré-requisito quanto com as disicplinas do semestre corrente, como SDCD e outras.
%Compreender, projetar e desenvolver sistemas robóticos móveis.
}

\def\PUDementa{Conceituação de robótica. Graus de liberdade. Atuadores e sensores em robótica industrial. Modelagem e simulação de robôs industriais. Programação de robôs industriais. Integração de robôs com periféricos no ambiente industrial. Introdução a Robótica Móvel; locomoção de robôs; Cinemática de robôs móveis;
%Introdução a Robótica Móvel;  locomoção de robôs;  Cinemática de robôs móveis;  percepção;  Visão de máquina aplicada a Robótica Móvel;  localização de robôs móveis;  planejamento e navegação;  exemplos de robôs autônomos;  aplicações.
}

\def\PUDconteudo{ & Conceito e definições
%Introdução à Robótica móvel. Conceitos básicos e aplicações. 
\\ 
 & Tipos de robôs e controladores 
 %Locomoção. Introdução;  Robótica móvel com pernas e com rodas. 
\\ 
 & Tipos de atuadores e sensores  
 %Cinemática em Robótica Móvel. Introdução;  restrições e modelos cinemáticos;  manobrabilidade;  espaço de trabalho e controle de movimento. 
\\ 
 & Métodos de ensino de robôs e de integração com sistemas de automação Unidade
 %Percepção. Sensores;  Visão Computacional aplicada à Robótica;  incerteza na representação e extração de atributos. 
\\ 
 & Simulação e programação \textit{off-line}
 %Localização. Introdução;  desafios da localização: ruído e aliasing;  localização baseada em navegação e soluções programadas;  representação de crença;  representação de mapas;  localização probabilística baseada em mapas;  sistemas de localização alternativos e construção autônoma de mapas.  
\\ 
 & Modelagem matemática e simulação 
 %Planejamento e navegação. Introdução;  competências para navegação: planejamento e reação.  Arquiteturas de navegação. 
\\ 
 & Integração com periféricos ; Introdução à Robótica móvel. conceitos básicos e aplicações. Locomoção. Introdução; Robótica móvel com rodas. 
 %Inteligência Computacional Aplicada à Robótica. Redes Neurais, Lógica Fuzzy, Algoritmos genéticos, classificadores aplicados à Robótica. 
 \\
 & Aplicações com tecnologias atuais
 }

\def\PUDmetodo{Aulas expositórias que podem ser teóricas e/ou práticas, onde as práticas no laboratório serão marcadas no decorrer da disciplina.}

\def\PUDavalia{A avaliação será feita com aplicação de provas teóricas e/ou práticas, além da possibilidade de inclusão de trabalhos e seminários no decorrer da disciplina.}

\def\PUDbibbasica{ & \href{www.biblioteca.ifce.edu.br/index.asp?codigo_sophia=62047}{Niku}, S. B.  Introdução à robótica - análise, controle, aplicações.  2º ed.  Rio de Janeiro, RJ: LTC, 2015.  382p.  ISBN: 9788521622376.
%\\ 
 %&  BEKEY, George A.  Autonomous robots: from biological inspiration to implementation and control.  Massachusetts (EUA): Massachusetts Institute of Technology - MIT, 2005.  577 p.   ISBN 9780262025782
\\ 
 & \href{www.biblioteca.ifce.edu.br/index.asp?codigo_sophia=23877}{Rosário}, João Maurício.  Princípios de mecatrônica.  São Paulo, SP: Prentice Hall, 2005.  356p.  ISBN: 9788576050100.
\\
 & \href{www.biblioteca.ifce.edu.br/index.asp?codigo_sophia=62087}{Affonso}, L. O. A.  Equipamentos mecânicos: análise de falhas e solução de problemas.  3º ed.  Rio de Janeiro, RJ: Qualitymark, 2014.  387p.  ISBN: 9788541400367. }

\def\PUDbibcomplementar{ & \href{www.biblioteca.ifce.edu.br/index.asp?codigo_sophia=41222}{Craig}, John J.  Robótica.  3º ed.  São Paulo, SP: Pearson, 2013.  379p.  ISBN: 9788581431284.
\\
 &  \href{www.biblioteca.ifce.edu.br/index.asp?codigo_sophia=65450}{Russell}, Stuart.  Inteligência artificial.  3º ed.  Rio de Janeiro, RJ : Elsevier, 2013.  988p.  ISBN: 9788535237016.
\\
 &  \href{www.biblioteca.ifce.edu.br/index.asp?codigo_sophia=24467}{Coppin}, Ben.  Inteligência artificial.  Rio de Janeiro, RJ : LTC, 2012.  636p.  ISBN: 9788521617297.
\\
 &  \href{www.biblioteca.ifce.edu.br/index.asp?codigo_sophia=46123}{Bräunl}, Thomas.  Embedded robotics: mobile robot design and applications with embedded systems.  3º ed.  Perth, Austrália: Springer, 2008.  541p.  ISBN: 9783540705338.
\\
 &  \href{www.biblioteca.ifce.edu.br/index.asp?codigo_sophia=56518}{Gonzalez}, Rafael C.; Woods, Richard E.  Processamento Digital de Imagens.  E-book.  Pearson.  644p.  ISBN: 9788576054016. } 
%&  MADRID, Marconi Kolm.  Curso sobre robôs industriais.  Fortaleza (CE): Universidade Federal do Ceará - UFC, 1992.  92 p.  

\def\PUDautor{Geraldo Luis Bezerra Ramalho}

\def\PUDdata{2013-04-17}

\def\PUDrevisao{3}

\def\PUDrevisor{Antonio Barbosa de Souza Júnior}

\def\PUDdatarevisao{2019-05-15}

\PUDcorpo
